\chapter{Breast Reconstruction}

\section*{Overview}
Breast reconstruction restores breast shape after mastectomy, lumpectomy, trauma, or congenital difference using an implant, the patient’s own tissue, or both.

\section*{Why This Test or Procedure Is Ordered}
It is considered to restore body contour and symmetry and may be performed immediately with mastectomy or later after cancer treatment and recovery.

\section*{How This Test or Procedure Is Performed}
Implant reconstruction may use a tissue expander followed by a permanent implant. Autologous reconstruction transfers tissue from areas such as the abdomen or back, sometimes using microsurgery to reconnect blood vessels.

\section*{Risks and Recovery}
Infection, bleeding, wound problems, capsular contracture, implant rupture, flap failure, asymmetry, and revision surgery are possible. Recovery includes wound care, activity restrictions, and staged nipple or symmetry procedures when desired.

\section*{References}
\begin{enumerate}
  \item MedlinePlus. Breast Reconstruction. U.S. National Library of Medicine; 2025.
  \item Current Medical Diagnosis and Treatment. McGraw-Hill; 2025.
\end{enumerate}
\clearpage
